\documentclass[addpoints]{exam}
% \documentclass[addpoints,12pt]{exam}

\usepackage[slovene]{babel}
\usepackage{amsfonts}
\usepackage{amssymb}
\usepackage{amsmath}
\usepackage{float}
\usepackage{graphicx}
\usepackage{enumitem}
\usepackage{color}
\usepackage{eurosym}


%Komentar naslednje vrstice glede na verzijo testa -- rešitve ali prostor za odgovore:
\printanswers

% \linespread{2}

\pointsinrightmargin
\marginbonuspointname{}


\renewcommand{\solutiontitle}{\noindent\textbf{Rešitve in točkovanje:}\par\noindent}

\colorgrids
\definecolor{GridColor}{gray}{0.7}
\setlength{\gridsize}{7mm}

\extrawidth{5mm}
\setlength{\rightpointsmargin}{1.5cm}

\begin{document}

\runningfooter{}{}{}

\begin{center}
    \fbox{\fbox{\parbox{15cm}{\centering
    Prvo pisno ocenjevanja znanja -- ponovno \\ 1. letnik -- test E \\ 6. november 2024 \\ (Osnove logike in teorije množic, Naravna in cela števila, Potence)}}}
\end{center}

\vspace{0.1in}
\makebox[\textwidth]{Ime in priimek, oddelek:\enspace\hrulefill}


\begin{center}
    \chqword{Naloga}
    \chpword{Možne točke}
    \chbpword{Dodatne točke}
    \chsword{Dosežene točke}
    \chtword{Skupaj}  
    \combinedpointtable[h][questions]
    % \combinedgradetable[h][questions]
\end{center}

\vspace{0.1in}


\begin{questions}

    % 1. vprašanje
    
    \question[8]
    Določite pravilnost izjave $(A\Rightarrow  B)\land \lnot C \Leftrightarrow (C\lor B)$ glede na izbor izjav $A$, $B$ in $C$. Upoštevajte vrstni red logičnih operacij.

    \begin{solution}[10cm]
        \begin{table}[H]
            \begin{tabular}{cccccccc}
             $A$ & $B$ & $C$ & $\lnot C$ & $A\Rightarrow B$ & $(A\Rightarrow B)\land\lnot C$ & $(C\lor B)$ & $(A\Rightarrow  B)\land \lnot C \Leftrightarrow (C\lor B)$\\ \hline
             P & P & P & N & P & N & P & N \\ 
             P & P & N & P & P & P & P & P \\
             P & N & P & N & N & N & P & N \\
             P & N & N & P & N & N & N & P \\
             N & P & P & N & P & N & P & N \\ 
             N & P & N & P & P & P & P & P \\
             N & N & P & N & P & N & P & N \\
             N & N & N & P & P & P & N & N
            \end{tabular}
        \end{table}
        Za vsako pravilno izpolnjeno vrstico, ki določa pravilnost, po $1$ točka.
    \end{solution}




    % 2. vprašanje
    \question[3]
    Določite resničnostno vrednost izjav:
    \begin{parts}
        \part 
        Ali velja $21\leq 21$ ali je to sestavljena izjava.\\
        \part 
        Če je $5>-6$, potem ni vsako naravno število večje od $0$.\\
        \part
        Ni res, da je $m(\mathbb{N})=\aleph_0$ natanko tedaj, ko je množica naravnih števil števno neskončna.\\
    \end{parts}

    \begin{solution}
        (a): P (prva izjava je pravilna, druga je pravilna) \\
        (b): N (prva izjava je pravilna, druga je nepravilna) \\
        (c): N (negirana izjava je nepravilna /prva izjava je pravilna, druga izjava je pravilna/) \\ \\
        Za vsako pravilno določeno logično vrednost po $1$ točka.
    \end{solution}



    \newpage

    % 3. vprašanje


    \question
    Množica $\mathcal{C}$ je prava podmnožica množice $\mathcal{B}$, prav tako je $\mathcal{A}$ prava podmnožica množice $\mathcal{B}$. Podmnožici $\mathcal{A}$ in $\mathcal{C}$ nimata praznega preseka.
    \begin{parts}
        \part[1]
        Grafično predstavite množico $(\mathcal{C}\cap\mathcal{A})\cup\mathcal{B}^\complement$.
        \part[2]
        Grafično predstavite množico $((\mathcal{C}\setminus\mathcal{A})\cup(\mathcal{A}\setminus\mathcal{C}))^\complement$.
    \end{parts}

    \begin{solution}[8cm]
        (a) in (b): grafična upodobitev \\
        Pri (a) in (b) za narisan Euler-Vennov diagram po $1$ točka.
    \end{solution}


%    \newpage 
    % 5. vprašanje
    \question
    V generaciji $99$ četrtošolcev so ugotovili, da so najbolj obiskane priprave na maturo iz treh izbirnih predmetov (\textit{geografije}, \textit{zgodovine} in \textit{psihologije}).
    Kar $67$ dijakov se pripravlja na \textit{geografijo}, $48$ na \textit{zgodovino} in $39$ na \textit{psihologijo}. Petintrideset dijakov obiskuje \textit{geografijo} in \textit{zgodovino},
    $16$ \textit{geografijo} in \textit{psihologijo} ter $10$ \textit{zgodovino} in \textit{psihologijo}. Dva dijaka se pripravljata na vse tri predmete.
    \begin{parts}
        \part[2]
        Koliko dijakov obiskuje priprave na maturo?
        \part[2]
        Koliko dijakov ne obiskuje priprav nobenega izmed predmetov?
    \end{parts}

    \begin{solution}[8cm]
        (a): $m(\mathcal{G}\cup\mathcal{Z}\cup\mathcal{P})=m(\mathcal{G})+m(\mathcal{Z})+m(\mathcal{P})-m(\mathcal{G}\cap\mathcal{Z})-m(\mathcal{G}\cap\mathcal{P})+m(\mathcal{Z}\cap\mathcal{P})+m(\mathcal{G}\cap\mathcal{Z}\cap\mathcal{P}) = 67+48+39-35-16-10+2=95$ \\
        (b): $m((\mathcal{G}\cup\mathcal{Z}\cup\mathcal{P})^\complement)=m(\mathcal{U})-m(\mathcal{G}\cup\mathcal{Z}\cup\mathcal{P})=99-95=4$ \\ \\
        Pri (a) in (b) za pravilen izračun $1$ točka, za zapis odgovora $1$ točka.
    \end{solution}



    % \newpage

    % 4. vprašanje
    \question
    Naj bo~ $\mathcal{U}=\left\{y;\ 4<y<21\right\}$. Dane so množice $\mathcal{A}=\left\{x;\ x=3n \land 2\leq n<6\right\}$,~ $\mathcal{B}=\left\{x;\ x=4k-1 \land 2\leq k\leq 5\right\}$ in $\mathcal{C}=\left\{x;\ x=5m \land 1<m\leq 4\right\}$. 
    \begin{parts}
        \part[5]
        Določite naslednje množice: $\mathcal{A}\cap\mathcal{C}\cap\mathcal{B}$, $\mathcal{C}\setminus\mathcal{A}$, $(\mathcal{A}\cup\mathcal{B})\setminus\mathcal{C}$ in $(\mathcal{A}\cup\mathcal{C})^\complement$.
        \part[2]
        Zapišite $\mathcal{PC}$. 
        \part[4]
        Izračunajte koliko elementov ima množica $\mathcal{PC}\times\mathcal{PB}$. Množico grafično predstavite, označite zgolj dva različna elementa.
    \end{parts}

    \begin{solution}[9cm]
        $\mathcal{A}=\left\{6,9,12,15\right\}$, $\mathcal{B}=\left\{7,11,15,19\right\}$, $\mathcal{C}=\left\{10,15,20\right\}$ \\
        (a): $\mathcal{A}\cap\mathcal{C}\cap\mathcal{B}=\left\{15\right\}$; $\mathcal{C}\setminus\mathcal{A}=\left\{10,20\right\}$; $(\mathcal{A}\cup\mathcal{B})\setminus\mathcal{C}=\left\{6,7,9,12,19\right\}$; $(\mathcal{A}\cup\mathcal{C})^\complement=\left\{5,7,8,11,13,14,16,17,18,19\right\}$ \\
        (b): $\mathcal{PC}=\left\{\emptyset, \{10\}, \{15\}, \{20\}, \{10,15\}, \{10,20\}, \{15,20\}, \{10,15,20\}\right\}$ \\
        (c): $m(\mathcal{PC}\times\mathcal{PB})=m(\mathcal{PC})\cdot m(\mathcal{PB})=2^{m(\mathcal{C})}\cdot 2^{m(\mathcal{B})} = 2^3\cdot 2^4=8\cdot 16 = 128$ + grafična upodobitev\\ \\
        Pri (a) za zapis/uporabo pravilnih množic $\mathcal{A}, \mathcal{B}, \mathcal{C}$ $1$ točka; za vsako pravilno določeno množico po $1$ točka. Pri (b) za zapis celotne $\mathcal{PC}$ $2$ točki, za zapis zgolj podmnožic $1$ točka. 
        Pri (c) za uporabo formule in izračun moči potenčne množice ter formule in izračun moči kartezičnega produkta po $1$ točka; za grafično upodobitev $1$ točka, za označena elementa $1$ točka.
    \end{solution}






    \newpage
        

    % 6. vprašanje
    \question[4]
    Izračunajte. $$ (-3)\left(1-3^2\right)^3+\left(3-(-2)^3+2\cdot(-6)\right)^{2n}$$

    \begin{solution}[7cm]
            $\begin{aligned}
                &(-3)\left(1-3^2\right)+\left(3-(-2)^3+2\cdot(-6)\right)^{2n}= \\
                    =& (-3)(1-9)^3+(3-(-8)-12)^{2n}=\\
                    =& (-3)(-512)+(-1)^{2n}= \\
                    =& 1536+1= \\
                    =& 1537
            \end{aligned} $ \\ \\
            Za upoštevanje sodih/lihih eksponentov glede na predznak $1$ točka, upoštevanje vrstnega reda operacij $1$ točka, izračun potenc $1$ in $-1$ $1$ točka, končen izračun $1$ točka.
    \end{solution}



    % 6. vprašanje
    \question[3]
    Poenostavite. $$ \left(-x^3\right)^2\cdot 10x^4\cdot\left(-x^5\right)^3$$

    \begin{solution}[7cm]
            $\begin{aligned}
                &\left(-x^3\right)^2\cdot 10x^4\cdot\left(-x^5\right)^3= \\
                    =& (x^3)^2\cdot 10x^4(-(x^5)^3)=\\
                    =& -x^6\cdot 10x^4\cdot x^{15}= \\
                    =& -10x^{6+4+15}= \\
                    =& -10x^{25}
            \end{aligned} $ \\ \\
            Za upoštevanje sodih/lihih eksponentov glede na predznak $1$ točka, izračun potenc z enako osnovo $1$ točka, končen izračun $1$ točka.
    \end{solution}


    
    %dodatna naloga
    \bonusquestion[3]
    \textit{Dodatna naloga:} Rešite enačbo $(X\setminus\{b\})\times(X \cap \{a,b\})=\left\{(a,a),(c,a),(d,a),(e,a)\right\}$.

    \begin{solution}
        $X=\left\{a,c,d,e\right\}$ \\
        Za v celoti pravilen zapis rešitve $3$ točke.
    \end{solution}


\end{questions}



\end{document}